\documentclass[12pt,a4paper]{article}

% ============================================================
% PAQUETES
% ============================================================

\usepackage[utf8]{inputenc}
\usepackage[T1]{fontenc}
\usepackage[spanish]{babel}
\usepackage{lmodern}
\usepackage{geometry}
\usepackage{microtype}
\usepackage{parskip}
\usepackage{enumitem}
\usepackage{tabularx}
\usepackage{array}
\usepackage{xcolor}
\usepackage{hyperref}

% ============================================================
% CONFIGURACIÓN
% ============================================================

\geometry{
    top=2.5cm,
    bottom=2.5cm,
    left=2.5cm,
    right=2.5cm
}

\setlength{\parindent}{0pt}

\hypersetup{
    colorlinks=true,
    linkcolor=black,
    urlcolor=blue
}

\setlist[itemize]{
    leftmargin=1.2cm,
    itemsep=0.2cm,
    topsep=0.2cm
}

\setlist[enumerate]{
    leftmargin=1.2cm,
    itemsep=0.2cm,
    topsep=0.2cm
}

\renewcommand{\arraystretch}{1.4}

% ============================================================
% DOCUMENTO
% ============================================================

\begin{document}

\begin{center}
    {\LARGE \textbf{Código de Ética, Convivencia y Trabajo Colaborativo}}\\[0.4cm]
    {\large Taller de Integración II y Taller de Integración IV}\\[0.3cm]
    Universidad Católica de Temuco
\end{center}

\vspace{0.5cm}

\section{Propósito}

El presente Código de Ética, Convivencia y Trabajo Colaborativo tiene como objetivo establecer criterios comunes de comportamiento, responsabilidad, comunicación y trabajo técnico para los integrantes de Taller de Integración II y Taller de Integración IV que participen conjuntamente en el desarrollo de un proyecto.

El documento busca facilitar el trabajo entre estudiantes pertenecientes a distintas generaciones académicas, establecer procedimientos conocidos previamente por todos los integrantes y evitar que los conflictos sean resueltos mediante criterios personales o improvisados.

Las disposiciones contenidas en este código deberán aplicarse considerando los principios de respeto, responsabilidad, transparencia, colaboración y trazabilidad del trabajo realizado.

\section{Principios generales}

Todos los integrantes deberán desarrollar sus actividades bajo los siguientes principios:

\begin{itemize}
    \item Respeto entre todos los integrantes del equipo.

    \item Responsabilidad individual respecto de las tareas, compromisos y funciones asumidas.

    \item Transparencia frente a errores, dificultades, retrasos o problemas que puedan afectar al proyecto.

    \item Colaboración entre los integrantes, especialmente cuando existan diferencias de conocimientos o experiencia.

    \item Trazabilidad de las decisiones y modificaciones relevantes realizadas durante el proyecto.

    \item Comunicación oportuna ante situaciones que puedan afectar el normal desarrollo del trabajo.
\end{itemize}

\section{Críticas y diferencias técnicas}

Se permite y fomenta la realización de críticas constructivas respecto de:

\begin{itemize}
    \item código;
    \item arquitectura;
    \item decisiones técnicas;
    \item metodologías;
    \item documentación;
    \item procedimientos de trabajo;
    \item organización del proyecto.
\end{itemize}

Toda crítica deberá estar orientada a mejorar el trabajo realizado.

No se considerará apropiado realizar críticas dirigidas directamente contra una persona sin explicar el problema identificado o sin plantear una alternativa, recomendación o posible solución.

Las discusiones deberán centrarse en el trabajo realizado y no en características personales del integrante responsable.

\section{Reuniones}

\subsection{Notificación}

Las reuniones deberán ser comunicadas con anticipación suficiente para permitir que los integrantes organicen adecuadamente su disponibilidad.

Se consideran medios oficiales de comunicación para efectos de notificación:

\begin{itemize}
    \item Discord;
    \item correo electrónico;
    \item Linear;
    \item WhatsApp.
\end{itemize}

Una reunión comunicada mediante cualquiera de estos medios se considerará oficialmente notificada.

\subsection{Disponibilidad de los integrantes}

Cuando una reunión esté siendo organizada, si un integrante informa inmediatamente que no puede participar en la fecha u horario propuesto, se deberá procurar encontrar una alternativa que permita la participación de todo el equipo.

Una vez establecida y notificada la reunión, un integrante que conozca con anticipación que no podrá participar deberá comunicarlo a más tardar un día antes de la realización de esta.

\subsection{Emergencias y situaciones de fuerza mayor}

Cuando exista una emergencia, imprevisto grave o situación de fuerza mayor que impida realizar una notificación con la anticipación requerida, el integrante deberá justificar posteriormente su ausencia ante el Scrum Master correspondiente.
\subsection{Ausencias}

Se considerará una ausencia justificada cuando:

\begin{itemize}
    \item haya sido informada dentro del plazo establecido;
    \item corresponda a una emergencia;
    \item exista una situación de fuerza mayor debidamente comunicada.
\end{itemize}

La ausencia a una reunión sin notificación previa ni justificación será considerada inicialmente una \textbf{falta leve}.

La persona que no participe en una reunión tendrá la responsabilidad de recuperar de manera autónoma la información tratada y ponerse al día respecto de:

\begin{itemize}
    \item temas discutidos;
    \item decisiones adoptadas;
    \item modificaciones realizadas;
    \item tareas asignadas;
    \item nuevos compromisos establecidos.
\end{itemize}

La ausencia de un integrante no obliga al resto del equipo a realizar nuevamente la reunión.

\subsection{Participación en decisiones}

Los integrantes ausentes deberán aceptar las decisiones adoptadas válidamente durante la reunión en la cual no participaron.

Al no encontrarse presentes durante la discusión y resolución correspondiente, no dispondrán de voto retroactivo respecto de dichas decisiones.

Esto no impedirá solicitar posteriormente una revisión cuando existan antecedentes técnicos, éticos o académicos relevantes que no hayan sido considerados originalmente.

\subsection{Puntualidad}

Al horario establecido para una reunión se otorgará un margen máximo de espera de \textbf{diez minutos} antes de iniciar formalmente.

Este margen se establece considerando que las reuniones tienen normalmente una duración aproximada de una hora y que retrasar su inicio afecta la disponibilidad del resto de los participantes.

Los atrasos ocasionales no constituirán automáticamente una falta. Sin embargo, la reiteración injustificada podrá ser considerada al evaluar el comportamiento del integrante.

\subsection{Registro de reuniones}

Toda reunión deberá dejar un registro mínimo que contenga:

\begin{itemize}
    \item fecha;
    \item hora de inicio;
    \item duración;
    \item temas tratados;
    \item decisiones tomadas, cuando corresponda.
\end{itemize}

Este registro tendrá como propósito mantener trazabilidad sobre el desarrollo del proyecto y permitir que integrantes ausentes puedan recuperar posteriormente la información.

\section{Gestión del repositorio GitHub}

\subsection{Principio general}

Todo desarrollo relacionado con el código fuente deberá realizarse mediante el uso de ramas.

Como regla general:

\begin{center}
    \textbf{No se permite desarrollar código directamente sobre la rama \texttt{main}.}
\end{center}

Cada funcionalidad, corrección o modificación deberá desarrollarse en una rama independiente y posteriormente integrarse mediante una Pull Request.

\subsection{Excepción para documentación}

Se permitirán commits directos a \texttt{main} únicamente cuando correspondan exclusivamente a:

\begin{itemize}
    \item documentación;
    \item archivos Markdown;
    \item diagramas;
    \item material explicativo;
    \item otros archivos de naturaleza documental.
\end{itemize}

Esta excepción será válida únicamente cuando la modificación no altere código fuente, configuración de ejecución, dependencias o comportamiento funcional del software.

\subsection{Ramas}

Cada tarea o modificación deberá utilizar una rama destinada específicamente a su desarrollo.

Las ramas creadas durante el proyecto constituyen parte de la documentación y de la evidencia del trabajo realizado por cada integrante. Permiten conservar trazabilidad respecto de las tareas desarrolladas, los cambios realizados y la participación de los distintos miembros del equipo.

Por este motivo:

\begin{center}
    \textbf{Queda prohibida la eliminación de ramas del repositorio GitHub.}
\end{center}

Esta prohibición aplica tanto a ramas que se encuentren actualmente en desarrollo como a ramas que ya hayan sido integradas mediante una Pull Request.

En consecuencia:

\begin{itemize}
    \item no se deberán eliminar manualmente ramas remotas;
    \item no se deberán eliminar ramas después de realizar un merge;
    \item deberá mantenerse desactivada la eliminación automática de ramas después de integrar una Pull Request;
    \item las ramas utilizadas deberán conservarse como evidencia y documentación del trabajo realizado;
    \item no se deberá eliminar trabajo perteneciente a otros integrantes;
    \item no se deberá modificar deliberadamente el historial para ocultar acciones realizadas;
    \item no se deberán eliminar commits relevantes con el objetivo de reducir o alterar la trazabilidad del proyecto.
\end{itemize}

Los movimientos relevantes realizados en el repositorio deberán conservar un nivel razonable de trazabilidad.

\subsection{Pull Requests}

Toda modificación de código deberá incorporarse a la rama principal mediante una Pull Request.

Antes de realizar el merge, la Pull Request deberá contar con al menos:

\begin{center}
    \textbf{una aprobación de un integrante de Taller de Integración IV.}
\end{center}

La aprobación representa una revisión adicional del trabajo, pero no elimina la responsabilidad del autor respecto de verificar previamente su implementación.

Asimismo, la existencia de una revisión implica una responsabilidad compartida sobre el proceso de integración. Si posteriormente se descubre un error que no era razonablemente detectable por el autor o que fue aprobado durante la revisión, no deberá utilizarse automáticamente como fundamento para responsabilizar exclusivamente a una persona.

\subsection{Método de integración}

Las Pull Requests deberán integrarse utilizando el mecanismo:

\begin{center}
    \textbf{Squash and Merge}
\end{center}

Este procedimiento permitirá consolidar los commits realizados durante el desarrollo de una tarea antes de incorporarlos a la rama principal, manteniendo un historial más ordenado y comprensible.

La realización del merge no implica la eliminación de la rama utilizada. La rama deberá permanecer disponible en GitHub como parte de la documentación y trazabilidad del proyecto.

\subsection{Force Push}

Queda prohibido utilizar mecanismos destinados a sobrescribir forzosamente el historial remoto, incluyendo:

\begin{verbatim}
git push --force
\end{verbatim}

y operaciones equivalentes.

Cuando exista una situación excepcional relacionada con el historial de Git, esta deberá ser discutida con el equipo antes de realizar cualquier acción que pueda afectar el trabajo remoto.

\subsection{Conflictos de merge}

Cuando una rama genere un conflicto al intentar integrarse con otra rama o con \texttt{main}, la responsabilidad inicial de resolver dicho conflicto corresponderá al integrante responsable de la rama que requiere ser integrada.

El responsable podrá solicitar ayuda a otros integrantes cuando sea necesario, pero no deberá trasladar automáticamente la resolución a otra persona.

\section{Contenido innecesario en el repositorio}

Se prohíbe incorporar al repositorio archivos, directorios o artefactos generados que no correspondan al código fuente o documentación que deba mantenerse versionada.

Ejemplos incluyen:

\begin{verbatim}
node_modules/
.cache/
dist/
build/
coverage/
.tmp/
\end{verbatim}

cuando estos elementos deban encontrarse correctamente definidos dentro de archivos como \texttt{.gitignore}.

La incorporación accidental de este tipo de material se considerará inicialmente una \textbf{falta leve}.

El integrante responsable deberá corregir el problema y realizar las modificaciones necesarias para evitar su repetición.

La reincidencia podrá aumentar la gravedad de la falta.

\section{Credenciales e información sensible}

Queda estrictamente prohibido almacenar dentro del repositorio:

\begin{itemize}
    \item archivos \texttt{.env} que contengan información real;
    \item contraseñas;
    \item tokens de acceso;
    \item claves de API;
    \item claves privadas;
    \item secretos de autenticación;
    \item credenciales de servicios externos;
    \item cualquier otra información cuya publicación pueda comprometer la seguridad del sistema.
\end{itemize}

Asimismo, queda prohibido incorporar este tipo de información directamente dentro del código fuente.

Por ejemplo:

\begin{verbatim}
const API_KEY = "clave-real-del-servicio";
\end{verbatim}

La incorporación de credenciales reales o secretos al repositorio será considerada una \textbf{falta grave}.

Cuando una credencial haya sido publicada accidentalmente, eliminarla posteriormente del archivo no será considerado suficiente, dado que puede permanecer disponible dentro del historial de Git.

La credencial deberá considerarse comprometida y deberán adoptarse las medidas correspondientes para:

\begin{itemize}
    \item revocarla;
    \item reemplazarla;
    \item modificar las credenciales relacionadas cuando corresponda;
    \item informar el incidente al equipo.
\end{itemize}

\section{Errores técnicos}

El equipo reconoce que los errores técnicos son una consecuencia natural del proceso de desarrollo de software.

Por este motivo:

\begin{center}
    \textbf{Cometer un error técnico no constituye por sí mismo una falta ética.}
\end{center}

La evaluación deberá considerar principalmente la conducta del integrante después de detectar o conocer el problema.

Por ejemplo:

\begin{itemize}
    \item Introducir accidentalmente un error y colaborar en su corrección no constituye una falta.

    \item Detectar un error propio e informar oportunamente al equipo constituye una conducta responsable.

    \item Detectar un problema relevante y ocultarlo deliberadamente podrá constituir una falta.

    \item Intentar eliminar evidencia con el objetivo de ocultar una acción podrá constituir una falta grave.

    \item Modificar intencionalmente el proyecto con el propósito de perjudicar su funcionamiento o el trabajo de otro integrante será considerado una falta muy grave.
\end{itemize}

\section{Criterios para determinar la gravedad de un incidente}

La gravedad de una situación no deberá determinarse exclusivamente por el resultado final producido.

Para evaluar un incidente deberán considerarse conjuntamente los siguientes elementos:

\begin{enumerate}
    \item \textbf{Intención:} si la acción fue accidental, negligente o deliberada.

    \item \textbf{Impacto:} nivel de afectación producido sobre el proyecto, equipo, usuarios o información.

    \item \textbf{Reincidencia:} existencia de incidentes similares anteriores.

    \item \textbf{Ocultamiento:} existencia de intentos por esconder el problema o evitar asumir responsabilidad.

    \item \textbf{Cooperación:} disposición del integrante para informar, corregir y prevenir nuevamente el problema.
\end{enumerate}

Por tanto, la evaluación general deberá considerar:

\begin{center}
\textbf{Intención + Impacto + Reincidencia + Ocultamiento + Cooperación}
\end{center}

\section{Clasificación de faltas}

Las situaciones podrán clasificarse en cuatro niveles generales.

\subsection{Falta leve}

Corresponde a incumplimientos menores, normalmente accidentales, cuyo impacto sobre el proyecto sea reducido.

Ejemplos:

\begin{itemize}
    \item ausencia injustificada ocasional a una reunión;
    \item incorporación accidental de archivos innecesarios al repositorio;
    \item actualización tardía de una tarea;
    \item incumplimiento menor de un procedimiento interno.
\end{itemize}

\subsection{Falta moderada}

Corresponde a incumplimientos reiterados o situaciones que comienzan a afectar directamente el trabajo de otros integrantes.

Ejemplos:

\begin{itemize}
    \item reiteración de faltas leves;
    \item ausencias injustificadas frecuentes;
    \item retrasos recurrentes sin comunicación;
    \item incumplimiento repetitivo de procedimientos establecidos;
    \item mantener tareas bloqueadas sin informar al equipo.
\end{itemize}

\subsection{Falta grave}

Corresponde a acciones que puedan comprometer significativamente el proyecto, su seguridad, su trazabilidad o el trabajo de otros integrantes.

Ejemplos:

\begin{itemize}
    \item incorporar credenciales reales al repositorio;
    \item realizar acciones destinadas a alterar de manera no autorizada el historial del proyecto;
    \item ocultar deliberadamente un incidente relevante;
    \item ignorar conscientemente procedimientos de seguridad;
    \item eliminar trabajo de otros integrantes sin autorización;
    \item eliminar deliberadamente ramas del repositorio y afectar con ello la trazabilidad del trabajo realizado.
\end{itemize}

\subsection{Falta muy grave}

Corresponde a conductas deliberadas destinadas a engañar, perjudicar, sabotear o apropiarse indebidamente del trabajo realizado.

Ejemplos:

\begin{itemize}
    \item sabotaje deliberado del proyecto;
    \item eliminación intencional de trabajo ajeno con intención de perjudicar;
    \item falsificación deliberada de aportes o trabajo realizado;
    \item apropiación del trabajo de otro integrante;
    \item publicación intencional de credenciales;
    \item acoso, hostigamiento o ataques personales graves.
\end{itemize}

\section{Reincidencia}

La clasificación inicial de una situación no impide que su gravedad aumente cuando exista reincidencia.

La repetición sistemática de faltas leves podrá ser evaluada como una falta moderada o, dependiendo de su impacto y circunstancias, como una falta de mayor gravedad.

Por ejemplo:

\begin{center}
\begin{tabularx}{0.9\textwidth}{
    >{\raggedright\arraybackslash}X
    >{\raggedright\arraybackslash}X
}
\textbf{Situación} & \textbf{Evaluación posible} \\ \hline

Una ausencia injustificada aislada
&
Falta leve
\\

Ausencias injustificadas recurrentes
&
Falta moderada
\\

Ausencias sistemáticas, abandono de tareas y ausencia de comunicación
&
Puede constituir una falta grave
\\

\end{tabularx}
\end{center}

\section{Gestión de tareas}

Las herramientas de gestión de tareas, especialmente Linear, deberán utilizarse como fuente formal para la asignación y seguimiento del trabajo.

\subsection{Triage y asignación de tareas de Taller de Integración II}

Los integrantes de Taller de Integración IV podrán crear, preparar, describir, priorizar y establecer los criterios de aceptación de las tareas correspondientes a Taller de Integración II.

Sin embargo, las tareas destinadas a Taller de Integración II deberán ser dejadas inicialmente en el estado:

\begin{center}
    \textbf{Triage}
\end{center}

Los integrantes de Taller de Integración IV no deberán asignar directamente estas tareas a integrantes específicos de Taller de Integración II.

La responsabilidad de revisar las tareas disponibles en Triage y distribuirlas dentro del equipo corresponderá exclusivamente al:

\begin{center}
    \textbf{Scrum Master de Taller de Integración II.}
\end{center}

En consecuencia, únicamente el Scrum Master de Taller de Integración II podrá:

\begin{itemize}
    \item retirar o mover una tarea de TI2 desde Triage para incorporarla formalmente al flujo de trabajo;
    \item asignar inicialmente una tarea a un integrante de TI2;
    \item cambiar al responsable de una tarea de TI2;
    \item reasignar una tarea de un integrante de TI2 a otro integrante;
    \item redistribuir tareas cuando existan problemas de carga de trabajo, bloqueos o dependencias.
\end{itemize}

Los demás integrantes de Taller de Integración II no podrán reasignarse tareas entre ellos ni modificar unilateralmente al responsable definido por el Scrum Master.

Por ejemplo, un integrante de TI2 que tenga una tarea asignada no podrá entregársela o asignársela directamente a otro compañero mediante Linear.

Cuando considere necesaria una reasignación, deberá comunicar la situación al Scrum Master de TI2, quien determinará si corresponde realizar el cambio.

Esta restricción se refiere exclusivamente a la \textbf{asignación y reasignación de responsables}. El integrante responsable de una tarea sí deberá mantener actualizado su estado de avance dentro del flujo de trabajo correspondiente, por ejemplo:

\begin{center}
\texttt{Todo} $\rightarrow$ \texttt{In Progress} $\rightarrow$ \texttt{In Review} $\rightarrow$ \texttt{Done}
\end{center}

De esta manera se mantiene una distribución coordinada del trabajo sin impedir que cada integrante refleje correctamente el progreso de sus propias tareas.

\subsection{Responsabilidad sobre las tareas asignadas}

Cada tarea deberá contar con un integrante responsable claramente identificado.

La persona asignada será responsable de:

\begin{itemize}
    \item conocer el objetivo de la tarea;
    \item mantener un seguimiento razonable de su estado;
    \item informar bloqueos;
    \item comunicar dificultades;
    \item informar retrasos;
    \item solicitar apoyo cuando sea necesario;
    \item evitar marcar como finalizado trabajo que aún no cumple sus requisitos.
\end{itemize}

Las herramientas de gestión de tareas, como Linear, deberán reflejar de manera razonable el estado real del trabajo.

Una tarea no deberá marcarse como completada únicamente para aparentar progreso.

\section{Tareas bloqueadas}

Cuando un integrante no pueda continuar una tarea debido a un bloqueo técnico, dependencia externa o falta de información, deberá comunicar la situación al equipo.

Solicitar ayuda frente a un problema técnico no constituye una falta ni será considerado una demostración de incapacidad.

Mantener deliberadamente una tarea detenida durante un periodo significativo sin comunicar el problema podrá ser considerado un incumplimiento de responsabilidad.

\section{Tareas atrasadas}

Cuando una tarea se atrase debido a circunstancias atribuibles principalmente al integrante responsable, este deberá:

\begin{itemize}
    \item informar el retraso;
    \item justificar la situación con su Scrum Master;
    \item reorganizar su planificación;
    \item disponer del tiempo adicional necesario para compensar razonablemente el atraso.
\end{itemize}

La responsabilidad de compensar un retraso propio no deberá trasladarse automáticamente a otros integrantes.

Cuando el retraso tenga su origen en:

\begin{itemize}
    \item dependencias externas;
    \item tareas previas no completadas por terceros;
    \item problemas técnicos ajenos al responsable;
    \item modificaciones de requerimientos;
    \item situaciones de fuerza mayor;
\end{itemize}

la situación deberá evaluarse considerando dichas circunstancias y no deberá atribuirse automáticamente como responsabilidad individual.

\section{Responsabilidad sobre el trabajo revisado}

El autor de una modificación es responsable de revisar y probar razonablemente su trabajo antes de solicitar su integración.

Sin embargo, las Pull Requests constituyen también un mecanismo de revisión colectiva.

Por ello, cuando una modificación haya sido revisada y aprobada conforme al procedimiento establecido, los errores posteriores deberán analizarse considerando tanto:

\begin{itemize}
    \item las responsabilidades del autor;
    \item el proceso de revisión;
    \item las pruebas existentes;
    \item las condiciones bajo las cuales ocurrió el error.
\end{itemize}

El objetivo de una revisión de código es aumentar la calidad del proyecto y no establecer posteriormente mecanismos de búsqueda individual de culpables.

\section{Procedimiento general frente a incidentes}

Cuando ocurra un incidente relevante, se procurará seguir el siguiente procedimiento:

\begin{enumerate}
    \item \textbf{Informar:} comunicar la situación a los integrantes correspondientes.

    \item \textbf{Contener:} evitar que el problema continúe generando consecuencias.

    \item \textbf{Registrar:} conservar información suficiente para comprender lo ocurrido cuando la situación lo requiera.

    \item \textbf{Analizar:} identificar la causa y alcance del problema.

    \item \textbf{Corregir:} aplicar una solución técnica u organizacional.

    \item \textbf{Evaluar:} determinar, cuando corresponda, la gravedad del incidente.

    \item \textbf{Prevenir:} establecer medidas que reduzcan la posibilidad de repetición.
\end{enumerate}

La prioridad inicial deberá ser resolver y contener el problema antes de determinar responsabilidades individuales.

\section{Medidas frente a incumplimientos}

Dependiendo de la naturaleza, gravedad y reincidencia de una situación, podrán considerarse medidas progresivas tales como:

\begin{itemize}
    \item conversación con el integrante;
    \item advertencia;
    \item registro del incidente;
    \item revisión de los procedimientos involucrados;
    \item corrección obligatoria del trabajo afectado;
    \item reasignación temporal de determinadas responsabilidades;
    \item comunicación al Scrum Master correspondiente;
    \item reunión entre los Scrum Masters de ambos talleres;
    \item comunicación al docente responsable cuando la gravedad o reincidencia lo justifique.
\end{itemize}

La aplicación de estas medidas deberá considerar siempre las circunstancias particulares del incidente y evitar sanciones desproporcionadas.

\section{Resolución de conflictos}

Cuando exista una diferencia entre integrantes, se deberá procurar inicialmente una solución directa mediante comunicación respetuosa.

Cuando el conflicto no pueda resolverse de esta manera, podrá solicitarse la intervención del Scrum Master correspondiente.

Si el conflicto afecta a integrantes pertenecientes a ambos talleres, los Scrum Masters de Taller de Integración II y Taller de Integración IV podrán intervenir conjuntamente.

Cuando la situación exceda las capacidades de resolución interna del equipo o corresponda a una conducta de especial gravedad, deberá ser comunicada al docente responsable.

\section{Responsabilidad colectiva}

El desarrollo del proyecto constituye un esfuerzo colectivo.

Las responsabilidades individuales deberán mantenerse claramente definidas, pero los problemas técnicos deberán abordarse procurando encontrar soluciones antes que culpables.

La responsabilidad individual adquiere especial relevancia cuando exista:

\begin{itemize}
    \item negligencia reiterada;
    \item ocultamiento;
    \item incumplimiento deliberado de procedimientos;
    \item falsificación;
    \item sabotaje;
    \item negativa injustificada a colaborar con la solución de un problema.
\end{itemize}

\section{Aceptación y revisión del código}

Todos los integrantes deberán conocer las disposiciones establecidas en este documento.

El código podrá ser modificado cuando la experiencia del proyecto demuestre la necesidad de incorporar nuevas reglas o modificar procedimientos existentes.

Toda modificación relevante deberá ser comunicada al equipo y quedar registrada.

\vfill

\begin{center}
\textbf{Taller de Integración II -- Taller de Integración IV}\\
Universidad Católica de Temuco
\end{center}

\end{document}
